Modelos generativos texto-imagem podem repetir padrões de representação presentes em seus dados de treinamento. Este trabalho desenvolve uma análise aplicada, quantitativa e exploratória de 64 imagens sintéticas de ambientes universitários, produzidas por quatro sistemas a partir de \textit{prompts} neutros e inclusivos. O BiasAuditFW organiza o processamento em duas unidades. DeepFace com RetinaFace produz predições no nível de cada rosto detectado, enquanto CLIP estima, uma vez por imagem, margens de diversidade racial percebida e de diversidade de apresentação de gênero. Os escores CLIP resultam da comparação entre protótipos textuais de diversidade e homogeneidade, com normalização L2 explícita. A detecção facial será confrontada com caixas anotadas por dois avaliadores. As predições de atributos também serão comparadas com FairFace como referência secundária, sem tratar os classificadores como fontes de identidade demográfica. Mann--Whitney e Wilcoxon serão aplicados separadamente às duas margens, sempre no nível da imagem e com correção de Holm. A interpretação do Efeito Patchwork dependerá da convergência entre a baseline humana, os escores semânticos e a composição facial.

\noindent\textbf{Palavras-chave:} análise de viés; IA generativa; visão computacional; CLIP; diversidade percebida; representação demográfica.
